\section{Discussion}\label{s:discussion}

The results obtained confirm that the simulator behaves correctly with the theoretical predictions of CHSH and E91 key distribution protocol. In particular, the observed violation of the classical bound \(S \leq 2\) in the secure channel experiments demonstrates that the simulator correctly reproduces non-classical correlations. 

The average CHSH value of \(2.44\) in the two-node experiment indicates a strong, though not maximal, violation of the CHSH inequality. This reduction from the theoretical maximum is expected due to noise introduced by channel losses over the 100\,km fiber link. 

In contrast, the insecure channel experiment shows a CHSH value below the CHSH threshold. This indicates that the entanglement correlations are being destroyed in the presence of Mallory. The detection and termination of the protocol in these cases demonstrates that CHSH violation is a reliable security indicator. 

The relationship between CHSH value and distance is consistent with the expected degradation of entanglement fidelity over distance. As distance increases, fidelity degradation occurs, which in turn reduces the CHSH value. The approximately linear decrease observed in the results aligns with expectations as the fidelity degradation is simulated as a constant of 0.14\%

The results show that higher dark count rates increase the simulated execution time. Additional dark count events trigger more detector cooldowns, reducing the efficiency of key generation and requiring more entangled-pair generation attempts to reach the target key length. Although the increase is moderate, the trend highlights the impact of detector noise on protocol performance.

Overall, the results show that the system successfully implements CHSH-based verification for E91 key distribution protocol and demonstrates that physical behaviour was simulated correctly.
